% sections/06-data-directory.tex — Data Directory and Input Handling

\section{Data Directory and Input Handling}

% ── Requirement ──
\subsection{Requirement}

Your project must read input data from a \filepath{data/} directory. This directory should contain multiple input files covering different test scenarios, edge cases, and performance benchmarks.

% ── Directory Structure ──
\subsection{Directory Structure}

\begin{lstlisting}[style=tree]
data/
|-- input_default.json ......... Default/primary input data
|-- input_small.yaml ........... Small input for quick testing
|-- input_large.csv ............ Large input for stress testing
|-- edge_case_empty.xml ........ Empty input edge case
|-- edge_case_single.toml ...... Single element edge case
|-- edge_case_max.txt .......... Maximum values edge case
|-- test_approach_1.yaml ....... Specific algorithm test
|-- test_stress.json ........... Stress test scenario
|-- test_invalid.json .......... Invalid input for error handling
|-- expected_output.json ....... Expected results for verification
+-- benchmark_large.bin ........ Binary data for performance tests
\end{lstlisting}

% ── File Naming Convention ──
\subsection{File Naming Convention}

\begin{center}
\begin{tabular}{l l}
\hline
\rowcolor{tableheader} \textbf{Prefix} & \textbf{Purpose} \\
\hline
\code{input\_*} & Regular input data \\
\code{edge\_case\_*} & Edge case testing \\
\code{test\_*} & Specific test scenarios \\
\code{benchmark\_*} & Performance benchmarking \\
\code{expected\_*} & Expected output for verification \\
\code{sample\_*} & Example/demo data \\
\hline
\end{tabular}
\end{center}

% ── Supported Formats ──
\subsection{Supported Formats}

Your project may use any of the following data formats:

\begin{itemize}
  \item \textbf{JSON} --- most common, well-supported in \cpp
  \item \textbf{YAML} --- human-readable, good for configuration
  \item \textbf{CSV} --- tabular data, simple to parse
  \item \textbf{XML} --- structured data with schema support
  \item \textbf{Plain text} --- simple key-value or line-based formats
  \item \textbf{Binary} --- for performance benchmarking with large datasets
\end{itemize}

\begin{notebox}[Format Flexibility]
You are free to choose any file format(s) for your test data. The examples in this document demonstrate various formats to show possibilities, but you may:
\begin{itemize}
  \item Use a single format (e.g., all JSON) for simplicity
  \item Mix formats based on what suits each test case
  \item Choose based on your parsing library preferences
\end{itemize}
Pick what works best for your project---consistency and clarity matter more than format variety.
\end{notebox}

% ── External Libraries ──
\subsection{External Libraries}

You may use external libraries for file parsing. All dependencies must be documented and installable via \makedeps.

\subsubsection{Recommended Libraries}

\begin{itemize}
  \item \textbf{JSON:}
    \begin{itemize}
      \item \href{https://github.com/nlohmann/json}{nlohmann/json} --- header-only, easy to use
      \item \href{https://github.com/Tencent/rapidjson}{RapidJSON} --- high performance
      \item \href{https://github.com/simdjson/simdjson}{simdjson} --- fastest JSON parser
    \end{itemize}
  \item \textbf{YAML:}
    \begin{itemize}
      \item \href{https://github.com/jbeder/yaml-cpp}{yaml-cpp}
    \end{itemize}
  \item \textbf{CSV:}
    \begin{itemize}
      \item \href{https://github.com/vincentlaucsb/csv-parser}{csv-parser}
      \item \href{https://github.com/ben-strasser/fast-cpp-csv-parser}{Fast C++ CSV Parser}
    \end{itemize}
\end{itemize}

% ── Dependency Installation ──
\subsection{Dependency Installation}

Your \makedeps{} target should handle downloading external libraries:

\begin{lstlisting}[style=makefile]
deps:
	@echo "Installing dependencies..."
	mkdir -p libs
	wget -q -O libs/json.hpp \
	  https://github.com/nlohmann/json/releases/latest/download/json.hpp
	# For system packages: sudo apt-get install libyaml-cpp-dev
\end{lstlisting}

For CMake projects, use \code{FetchContent}:

\begin{lstlisting}[style=cpp]
# CMakeLists.txt
include(FetchContent)

FetchContent_Declare(
    json
    URL https://github.com/nlohmann/json/releases/latest/download/json.hpp
    DOWNLOAD_NO_EXTRACT TRUE
)
FetchContent_MakeAvailable(json)

target_link_libraries(${PROJECT_NAME} PRIVATE nlohmann_json::nlohmann_json)
\end{lstlisting}

% ── Input Handling Best Practices ──
\subsection{Input Handling Best Practices}

\begin{itemize}
  \item \textbf{Validate all input} before processing
  \item \textbf{Handle missing or malformed files} gracefully with clear error messages
  \item \textbf{Support command-line arguments} for specifying input files
  \item \textbf{Log errors} with context (file name, line number, expected format)
\end{itemize}

\begin{lstlisting}[style=cpp]
int main(int argc, char* argv[]) {
    std::string inputFile = "data/input_default.json";

    if (argc > 1) {
        inputFile = argv[1];
    }

    // ... load and process inputFile
}
\end{lstlisting}

% ── Test Data Examples ──
\subsection{Test Data Examples}

Below are brief structural examples. See \textbf{Appendix A} for extended themed examples.

\subsubsection{Basic Structure (JSON)}

\begin{lstlisting}[style=json]
{
    "nodes": [
        {"id": 1, "label": "A"},
        {"id": 2, "label": "B"},
        {"id": 3, "label": "C"}
    ],
    "edges": [
        {"from": 1, "to": 2, "weight": 5},
        {"from": 2, "to": 3, "weight": 3},
        {"from": 1, "to": 3, "weight": 10}
    ]
}
\end{lstlisting}

\subsubsection{Edge Cases}

Empty input:

\begin{lstlisting}[style=json]
{
    "nodes": [],
    "edges": []
}
\end{lstlisting}

Single node:

\begin{lstlisting}[style=json]
{
    "nodes": [{"id": 1, "label": "A"}],
    "edges": []
}
\end{lstlisting}

Maximum values:

\begin{lstlisting}[style=json]
{
    "nodes": [
        {"id": 2147483647, "label": "MAX"}
    ],
    "edges": [
        {"from": 2147483647, "to": 2147483647, "weight": 2147483647}
    ]
}
\end{lstlisting}

\subsubsection{Expected Output}

\begin{lstlisting}[style=json]
{
    "input_file": "data/input_default.json",
    "expected": {
        "shortest_path": ["A", "B", "C"],
        "total_cost": 8
    }
}
\end{lstlisting}

See \textbf{Appendix A} for extended themed examples.
